% Part: lambda-calculus
% Chapter: introduction
% Section: church-rosser

\documentclass[../../../include/open-logic-section]{subfiles}

\begin{document}

\olfileid{lam}{int}{cr}
\olsection{Sipat Church--Rosser}

\begin{thm}
\ollabel{thm:church-rosser}
Upamana $M$, $N_1$, lan $N_2$ iku suku
saengga $M \red N_1$ lan $M \red N_2$.
Mula ana suku $P$ saengga $N_1 \red P$
lan $N_2 \red P$.
\end{thm}

\begin{cor}
Upamana $M$ bisa direduksi dadi wujud normal.
Mula wujud normal iki unik.
\end{cor}

\begin{proof}
Yen $M \red N_1$ lan $M \red N_2$, miturut
teorema sadurunge ana suku~$P$ saengga
$N_1$ lan $N_2$ loro-lorone direduksi dadi~$P$.
Yen $N_1$ lan~$N_2$ loro-lorone wis wujud normal,
iki mung bisa kedadeyan yen
$N_1 \ident P \ident N_2$.
\end{proof}

Pungkasan, rong suku $M$ lan~$N$ diarani
\emph{ekuivalen-$\beta$}, utawa mung
\emph{ekuivalen}, yen loro-lorone bisa
direduksi dadi suku kang padha; yaiku
yen ana $P$ saengga $M \red P$ lan
$N \red P$. Iki ditulis $M \equal[\beta] N$.
Nganggo \olref{thm:church-rosser}, bisa
dipriksa yen $\equal[\beta]$ iku relasi
ekuivalensi, kanthi sipat tambahan:
kanggo saben $M$ lan~$N$, yen $M \red N$
utawa $N \red M$, mula
$M \equal[\beta] N$. (Malah,
bisa dibuktekake yen $\equal[\beta]$
iku relasi ekuivalensi \emph{paling cilik}
kang nduweni sipat iki.)

\end{document}
